\section{From the collinear stress identity to global Ward identities}
\label{sec:physical-ward-propagation}

This section records the finite-configuration argument which turns a collinear
rotation identity into a Ward identity at every separated configuration.  It
also records the two finite-level bridges needed later: the outer stress
estimate supplies the collinear identity, and erasure of the internal Ward
kernel turns the differentiated Ward identity into the Lie relations.  All
statements are at fixed finite insertion number and fixed field grades; no
coincidence limit is taken.

Let $\mathscr S$ be a Euclidean hierarchy.  Write $\mathcal F_d$ for its
complex field space at grade $d$, $\mathcal P_d:\mathcal F_d\to
\mathcal F_{d+1}$ for the positional derivative, and
$\mathcal Y_d(K):\mathcal F_d\to\mathcal F_d$ for a spin-row family indexed
real-linearly by a skew matrix $K\in\mathfrak{so}(4)$.  A label at a slot is
a pair $(d,v)$ with $v\in\mathcal F_d$.  For a word $q=(q_1,\ldots,q_n)$ and
a separated tuple $x=(x_1,\ldots,x_n)$, let $S(q;x)$ denote its Euclidean
moment.  The hierarchy's derivative rule says that differentiation in the
coordinate $b$ of $x_i$ replaces $q_i=(d_i,v_i)$ by
$(d_i+1,\mathcal P_{d_i,b}v_i)$.

Define the angular Ward expression by
\begin{align}
\mathcal A_K(q;x) = \sum_{i=1}^n \bigg(&S(q[i\mapsto
  (d_i,\mathcal Y_{d_i}(K)v_i)];x) \\
  &+\sum_{b=1}^4 (Kx_i)_b\,
  S(q[i\mapsto(d_i+1,\mathcal P_{d_i,b}v_i)];x)\bigg).
\label{eq:ward-angular-expression}
\end{align}
Here $q[i\mapsto\cdot]$ changes only the indicated slot.  The condition
\emph{Still at $x$} means $\mathcal A_K(q;x)=0$ for every word $q$ and every
$K$.  The condition \emph{collinear} means this identity holds whenever the
points lie on one common coordinate line and remain pairwise distinct.
The compatibility relation between position and spin derivatives is
\begin{equation}
 \mathcal P_{d,b}\mathcal Y_d(K)v
 +\sum_{c=1}^4 K_{cb}\,\mathcal P_{d,c}v
 =\mathcal Y_{d+1}(K)\mathcal P_{d,b}v.
\label{eq:ward-commutation}
\end{equation}
The signs in this identity use the convention in
\eqref{eq:ward-angular-expression}.

\begin{theorem}[Propagation from a line]\label{thm:ward-global-propagation}
Assume the hierarchy is smooth on the configuration space
\[
 \operatorname{Conf}_n=\{(x_1,\ldots,x_n):x_i\ne x_j\text{ for }i\ne j\},
\]
the spin rows obey \eqref{eq:ward-commutation}, and the angular expression
\eqref{eq:ward-angular-expression} is zero at every separated collinear
configuration.  Assume also the extreme-segment analyticity established in
Section~\ref{sec:physical-ward-analytic-continuation}: on every coordinate
segment used below, the angular expression for each fixed word and generator
is real analytic on a neighborhood of the segment.  Then it is zero at every
$x\in\operatorname{Conf}_n$.
\end{theorem}

\begin{proof}
We first establish a geometric reduction.  Fix a coordinate axis $a$ and
another axis $b\ne a$.  Call a tuple $a$-ordered if the real numbers
$x_i\cdot e_a$ are pairwise distinct.  Such tuples are dense in
$\operatorname{Conf}_n$: perturb the $a$-coordinate of point $i$ by
$t i$.  For each pair $i\ne j$, equality after perturbation can occur for at
most one value of $t$.  Since there are finitely many pairs, there are
arbitrarily small $t$ avoiding all these exceptional values.  The perturbed
tuple converges to the original one and remains in the open set
$\operatorname{Conf}_n$ for small enough $t$.

Every $a$-ordered tuple can be moved to a collinear tuple by a finite chain
of single-point coordinate segments which stay in $\operatorname{Conf}_n$.
Here are the construction and the separation check.  Let $I$ be the set of
points not yet on the $a$-axis, maintaining the invariant that points
already aligned lie beyond every remaining point in the $a$-coordinate.
Initially $I$ is the whole set, so the invariant is vacuous.  Choose
$i\in I$ with smallest $a$-coordinate.  First move point $i$ in coordinate $b$ far enough in the
positive direction that it is the unique largest $b$-coordinate.  During
this first segment it stays the unique smallest $a$-coordinate, since its
$a$-coordinate is fixed.  Then move it in coordinate $a$ far enough to be
the unique largest $a$-coordinate.  During this second segment it stays
the unique largest $b$-coordinate, so its $b$-coordinate does not meet
another point's.  Now, one at a time, set each of its other coordinates
to zero.  These moves leave its $a$-coordinate fixed and it remains the
unique largest in $a$.  Every intermediate configuration is therefore
separated.  The point has now reached the $a$-axis without moving any other
point, and it lies beyond every point still in $I$.  Repeat with the
remaining finite set $I$.  This terminates after at most $n$ stages and
gives the required finite chain.  The motion consists of finitely many
segments; the segment parameter interval can be oriented in either direction.

It remains to show the Ward identity is constant along each such segment.
Fix a segment moving point $i$ in coordinate $b$, with parameter $s$ in a
closed interval $[u,v]$.  By construction, throughout this interval point
$i$ is either the strict minimum or the strict maximum in some coordinate
$a$; the choice may switch between segments, but not within this one.  For
each fixed word $q$ and skew matrix $K$, the scalar function
\[
 s\longmapsto \mathcal A_K(q;x(s))
\]
is real analytic on a neighborhood of the interval by the stated
extreme-segment analyticity input.  This input applies to moment rows with
the moving point at either end and to every raised row appearing in the
angular expression.  It is precisely the result of
Section~\ref{sec:physical-ward-analytic-continuation}; smoothness on
$\operatorname{Conf}_n$ alone is not enough.

At an endpoint where Still holds, all derivatives of this analytic scalar
function vanish.  For order zero this is the assumed Ward identity.  For
higher orders, differentiate the identity along the segment and use
\eqref{eq:ward-commutation}: differentiating the spin substitution at the
moving slot gives the spin substitution of the raised positional derivative;
the derivative of the orbital coefficient supplies the matrix term
$\sum_c K_{cb}\mathcal P_{d,c}$; these terms combine exactly into the
raised-slot term in the angular identity.  Induction therefore identifies
the $r$th derivative with the Ward expression for the word with that slot
raised $r$ times.  It is zero at the endpoint.  A real-analytic function on
a connected interval whose derivatives of every order vanish at one point
vanishes on the interval: its convergent Taylor series is zero near that
point, and the identity theorem propagates the equality across the interval.
Thus Still transfers from the final endpoint of a segment to its initial
endpoint.  Applying this successively along the finite chain proves Still
at every $a$-ordered tuple.  Finally, each angular expression is continuous
on $\operatorname{Conf}_n$ by smoothness of the hierarchy there.  Density of
the ordered tuples proves it vanishes at the original arbitrary separated
tuple.  This also covers $n=0,1$, where there are no pairwise collision
constraints and the sums are empty or single-term.
\end{proof}

\paragraph{The collinear identity from the stress tensor.}
The preceding theorem needs no stress tensor once its collinear hypothesis
is available.  In the application, let $J_{ab}$ be a symmetric conserved
stress insertion, let $p_{d,b}$ be the local positional response, and let
$Y_d(K)$ be the spin response.  The stress-law package consists of the
conservation law, the symmetric raw stress identity, and the flux identity
that identifies the infinitesimal rotation of a moment with the integral
of the stress flux.  The actual source constructs the flux test rather
than assuming a noncompact rotation field is admissible.

For fixed separated collinear $x$, choose smooth compactly supported
windows $\chi_i$ around each point, equal to one near that point and with
supports avoiding all other insertions.  For a plane rotation in axes
$a,b$, the packet
\[
 F_c(y)=\sum_i\chi_i(y)(K y)_c
\]
is smooth, compactly supported, and agrees with the rotational vector
field near each insertion.  The flux identity therefore evaluates the
angular expression as minus the integral of its stress-flux density.  In
that density the term $\chi_i\,\partial_a(Ky)_b$ cancels after contraction
against the symmetric stress, because $K$ is skew.  The terms
$(\partial_a\chi_i)(Ky)_b$ do not cancel: they are the shell integrand.
Their supports lie in the cutoff shells around the insertions, and the
exterior packet estimate below controls them.

Here is the rate calculation behind the stress-packet estimate.  For a
sequence of scalar or finite-dimensional vector quantities $F_n$, write
$F=O_r$ when $\|F_n\|\leq C e^{rn}$ for all sufficiently large $n$, with
$C$ independent of $n$.  The exact choice of the fixed axes, word, and
positions may affect $C$.  A positive spectral pairing satisfies the
following elementary interpolation rule: if $X_n,Y_n\geq0$ and
\[
 |F_n(p)|^2\leq (M e^{tn})^2X_n(p)Y_n(p),\qquad
 \int X_n\,d\nu_B=O_x,\qquad \int Y_n\,d\nu_A=O_y,
\]
then Cauchy--Schwarz gives
\begin{equation}
 |\langle F_n\rangle_{B,A}|
 \leq M(C_XC_Y)^{1/2}e^{(t+(x+y)/2)n},
 \qquad\text{so }\quad
 \langle F_n\rangle_{B,A}=O_{t+(x+y)/2}.
\label{eq:ward-balanced-price}
\end{equation}
Here $C_X,C_Y$ are the constants in the respective integral bounds.  This
rule applies to the cross-spectral pairing after its positive supports
and integrability have been verified.

For the packet under consideration the anisotropic scale exponents are
\[
 w_j=\begin{cases}
 -100,&j=a\text{ or }j=b,\\
 -150,&j=l\text{ and }j\notin\{a,b\},\\
 -300,&\text{otherwise},
 \end{cases}
 \qquad
 u=\begin{cases}1,&d=l,\\-w_l,&d\ne l.\end{cases}
\]
They obey $-300\leq w_j\leq-100$, $w_d<0$, and $u\leq150$.  The
nonnegative spectral measures are supported on $p_d\geq0$.  The cove is
the set where $p_d\leq1$ and every transverse coordinate satisfies
$|p_j|\leq e^{(1+w_j)n}$.  Let $h$ be the fixed nonnegative, smooth, even
one-dimensional cutoff which equals one on $[-1,1]$ and is supported in
$[-2,2]$.  The transverse panel uses
the normalized dilation
\[
 h_{j,n}(s)=e^{-(1+w_j)n}h\bigl(s e^{-(1+w_j)n}\bigr),
 \qquad H_n(p)=\prod_{j\ne d}h_{j,n}(p_j).
\]
Thus $H_n$ includes the normalization factors of the three transverse
one-dimensional dilations; they are part of both price weights below.
The packet estimate is for pairwise distinct axes $a,b$ and for
$d\ne a$, as in the stress-packet construction.  Denote by $D_{r,n}$ the
positive longitudinal cost
\[
 D_{r,n}(p)=\mathbf1_{p_d\leq1}\,(p_d+e^{rn})^{-1}
             H_n(p),
\]
Denote by $Z_{r,n}$ its cone-weighted variant
\[
 Z_{r,n}(p)=\mathbf1_{p_d\leq1}
     \frac{p_d e^{4rn}}{(p_d+e^{rn})^4}H_n(p).
\]
The following are the precise price bounds needed from the collinear
spectral estimates.  The positive line and point prices are supplied by
Section~\ref{sec:physical-ward-packet-prices}; the cone-weighted bound is
proved by the fixed-band stress-conservation argument in
Section~\ref{sec:physical-ward-stress-cone}.  In these notations the
required bounds are
\begin{align}
 \int D_{-3000,n}\,d\nu_B&=O_u,&
 \int Z_{-99,n}\,d\nu_A&=O_{-395},
 \label{eq:ward-deep-prices}\\
 \int D_{-3000,n}\,d\nu_B&=O_u,&
 \int D_{0,n}\,d\nu_A&=O_{900},
 \label{eq:ward-shallow-prices}\\
 \int D_{-3000,n}\,d\nu_B&=O_u,&
 \int D_{w_d,n}\,d\nu_A&=O_1.
 \label{eq:ward-side-prices}
\end{align}
Each bound is uniform in $n$ for a fixed packet.  The exterior Fourier
estimate supplies $O_{-101}$ on the complement of the low cove by the
damped exterior estimate, not by assuming an undamped spectral moment.  The
distinction between the two price measures in each display is essential:
the first controls the collinear line cost, while the second controls the
packet amplitude.

The one-dimensional panel profile $\varphi_d$ in the test has Fourier gap
$\widehat\varphi_d(\xi)\ne0\Rightarrow |\xi|\geq1$.  This is a property
of the panel profile, not a claim that the compactly supported spatial
petal has a Fourier gap.  On the cove, the frozen Fourier estimate for that panel gives
the differentiated-packet bound
\begin{equation}
 |F_n(p)|\leq M e^{3n}H_n(p)
                    \exp(-p_de^{-w_dn}),\qquad p_d\geq0.
\label{eq:ward-packet-tail}
\end{equation}
On the cove with $d=b$, split at the longitudinal level
$p_d=e^{-3000n}$.  On the region $p_d\geq e^{-3000n}$, the elementary inequality
\[
 e^{-t/s}\leq 4\left(\frac{s}{t+s}\right)^2
 \quad(t\geq0, 0<s\leq v)
\]
with $s=e^{w_dn}$ and $v=e^{-99n}$ first gives the square-tail factor
$16(v/(p_d+v))^4$.  On this region, put $R=e^{-3000n}$.  Since
$R\leq p_d\leq1$, we have $1\leq4p_d/(p_d+R)$, and therefore
\[
 |F_n|^2\leq (M' e^{3n})^2D_{-3000,n}Z_{-99,n}.
\]
The factor $4p_d/(p_d+R)$ pays for replacing the packet's denominator at
scale $e^{-99n}$ by the longitudinal cost at scale $e^{-3000n}$.  The costs in
\eqref{eq:ward-deep-prices} and
\eqref{eq:ward-balanced-price} give exponent
\[
 3+\frac{u-395}{2}\leq 3+\frac{150-395}{2}=-119.5.
\]
On the other cove region, $0\leq p_d<e^{-3000n}$, the exponential in
\eqref{eq:ward-packet-tail} is at most one.  The packet square is therefore
bounded with amplitude exponent $3-1500=-1497$ against
$D_{-3000,n}D_{0,n}$.  The shallow prices yield
\[
 -1497+\frac{u+900}{2}\leq -972.
\]
Both estimates are stronger than $O_{-101}$.

If $d\ne b$, the shear estimate bounds the squared packet by an amplitude
with exponent $3+w_d$, against $D_{-3000,n}D_{w_d,n}$.  When $d=l$,
$w_d=-150$ and $u=1$, so \eqref{eq:ward-balanced-price} gives
$-150+3+(1+1)/2=-146$.  When $d\ne l$, $w_d=-300$ and $u\leq150$,
so the exponent is at most
\[
 3-300+\frac{150+1}{2}=-221.5.
\]
The cove is partitioned by these two cases, and the complement has the
independent damped Fourier-tail price $O_{-101}$.  For clarity, outside the
cove there are exactly two alternatives.  If $p_d>1$, the gap in the panel
Fourier support and the positive-time damping estimate give
\[
 e^{-p_de^{-w_dn}}\leq e^{yn}e^{-Rp_d}
\]
eventually, for any fixed $R>0$ and prescribed $y$.  Take $y=-1001$;
the transverse readout contributes at most $e^{900n}$, leaving rate
$-101$.  The positive-time damped spectral moment is supplied by the
exterior-price estimate; no undamped zeroth moment is assumed.  If
instead $p_d\leq1$
and a transverse scatter bound fails, then for some $j\ne d$,
$|p_j|>e^{(1+w_j)n}$.  The Schwartz bound of order $k$ in that coordinate
contributes the factor $e^{-kn}$, while the order-zero bounds for the two
other transverse coordinates and the ordinary bound in the failed
coordinate contribute at most $e^{900n}$.  Choose $k>1001$; the
corresponding positive-time polynomially weighted estimate is again the
exterior-price input.  These two cases give the unamplified matching
cluster readout the required price $O_{-101}$ under that input.

The cross-slice step needs absolute integrability on both sides.  The
weighted line/exterior prices and the positive-time damped spectral moment
first give the spectral-side integrability.  On the position side, place the
two fixed clusters at a common positive-time separation $u$.  The panel's
temporal Fourier gap is at least $e^{-w_dn}$, which is eventually greater
than $2u$ with a strict margin.  The damped moment estimate, together with
the Schwartz bounds in the transverse coordinates, then gives absolute
integrability and justifies Fubini.  This is the gap estimate in the
cross-slice input; weighted spectral Cauchy by itself would not prove
position-side integrability.  With these integrability hypotheses in place,
the cross-slice identity gives, for a fixed sign
$\epsilon\in\{-1,1\}$ and fixed $\tau>0$,
\[
 \epsilon\int R_n(p)\,dp=\tau\,\operatorname{cross}(F_n),
\]
where $R_n$ is the unamplified cluster readout times the position test and
$F_n$ is its slice packet.  Thus the cove and exterior estimates imply
$\int R_n=O_{-101}$.
The actual insertion packet has an additional exact scalar factor
$e^{100n}$: the profile identity is
$\mathrm{packet}_n=e^{100n}\,\mathrm{readout}_n$.  Consequently its
integral has price
\begin{equation}
 |B_{a,b,l,d}(n)|\leq C e^{100n}e^{-101n}=C e^{-n}.
\label{eq:ward-shell-budget}
\end{equation}
To pass from this packet price to the stress insertion, first represent the chosen
stress component by its finite signed cluster expansion.  Decompose the
other insertion labels into their finitely many sign patterns.  A coordinate
reflection changes the sign of each nonmatching pattern while preserving
measure, so its integral is zero.  For the matching pattern, the exact
derivative-profile identity is
$\partial_d\mathrm{petal}_n=e^{100n}\mathrm{positionTest}_n$; thus its
unamplified readout has price $O_{-101}$ and the actual stress packet has
price $O_{100-101}=O_{-1}$.  Finally, the finite sum over patterns, axes,
and insertions preserves $O_{-1}$, and the shell flux tends to zero.

For the plane with generator normalized by $K e_a=e_b$ and
$K e_b=-e_a$, use the compact
vector test with components
\[
 F_b^{(n)}(y)=\mathrm{petal}(a,b,l,n)(y),\qquad
 F_a^{(n)}(y)=-\mathrm{petal}(b,a,l,n)(y),\qquad
 F_c^{(n)}(y)=0\quad(c\notin\{a,b\}).
\]
The local packet limits are $F_b^{(n)}(y)\to y_a$ and
$F_a^{(n)}(y)\to-y_b$, exactly the two components of the orbital field.
The second component carries the minus sign; it is not an additional sign
in the stress contraction.  By stress symmetry, its shell integral and the
first component combine as
\[
 \int\mathrm{river}(F^{(n)})
 =\int\mathrm{draught}_{a,b}^{(n)}
  -\int\mathrm{draught}_{b,a}^{(n)}.
\]
Each draught is a finite sum of the already estimated shell terms, hence
has bound $C e^{-n}$.  The compact-flux identity identifies the negative of
this integral with the local plane-rotation response.  That response is
independent of $n$, so its absolute value is at most $C e^{-n}$ for all
sufficiently large $n$ and therefore is zero.  The generators of coordinate
planes span $\mathfrak{so}(4)$, giving the collinear identity for every
skew generator.
The skew matrices
$E_{ab}-E_{ba}$ form a real basis of $\mathfrak{so}(4)$, and the angular
response is real-linear in its generator, so it vanishes for every $K$.
This proves the collinear input used above.  The exponential estimate is
the substantive analytic input here; it is not a consequence of the
geometric propagation argument.

\begin{lemma}[Outer translation identity]\label{lem:ward-outer-miles}
Under the same stress-law assumptions, the vector positional response also
satisfies the following line identity.  If $x_1,\ldots,x_n$ lie on the
axis $b$, $j\ne b$, and $v_i\in\mathcal F_d$, then
\[
 \sum_{i=1}^n S\bigl(q[i\mapsto(d+1,p_{d,j}v_i)];x\bigr)=0,
 \qquad q_i=(d,v_i).
\]
\end{lemma}

\begin{proof}
Let $K_{j,b}$ be the plane generator normalized by $K_{j,b}e_b=e_j$.
Simultaneous translation preserves the moment and all spin terms.  In the
orbital term, translating every point by $e_b$ changes $(K_{j,b}x_i)_c$
to $(K_{j,b}x_i)_c+(K_{j,b}e_b)_c$.  Only the $c=j$ component changes,
by one.  Therefore
\[
 \mathcal A_{K_{j,b}}(q;x+e_b)=\mathcal A_{K_{j,b}}(q;x)
   +\sum_i S(q[i\mapsto(d+1,p_{d,j}v_i)];x).
\]
Both tuples remain separated and lie on the same line, so both angular
terms vanish.  Subtracting gives the stated line identity.  For the empty
word it is the empty sum.
\end{proof}

\paragraph{Finite-level bridges used by the selected scale.}
The line identity and the collinear rotational identity are only the
geometric part of the finite-level argument.  The selected scale uses the
following additional bridges; their order matters.

At the local algebraic stage, the three identities used above are derived
from the stress-law package rather than postulated separately.  For the
positional response, differentiate the compactly supported flux with the
same cutoff kept fixed.  The derivative has two terms: differentiation of
the flux coefficient, and differentiation of the moment at the moving
insertion.  The latter raises exactly that slot by one derivative.  The
stress conservation identity moves derivatives of the cutoff to the
coefficient, and symmetry of $J$ cancels the antisymmetric part.  This
gives the shell-commutation relation and compatibility of $p$ with grade
inclusion.  Applying translations, reflections, and real conjugation to
the same cutoff formula gives the color covariance, reality, unit, and
mirror rules for $p$ and $Y$.  In particular the unit rule uses
conservation on the word with the unit insertion deleted; it is not a
formal consequence of a nonzero-grade relation.

To obtain the all-order line jets, do not differentiate an identity known
only on the line.  First form the mixed-cap row which sums, over the marked
slot, the moment with that label raised once in a fixed transverse direction.
The derivative rule and the positional/spin compatibility identity prove
on the whole open configuration space that differentiating this row in any
point-coordinate direction equals the corresponding row with that slot
raised once more.  This is a local identity on $\operatorname{Conf}_n$,
not a line-restricted assertion.  On the line, the order-zero mixed-cap row
vanishes: pad all caps to their finite maximum by inclusion compatibility,
then apply the line-sum identity.  Now induct on derivative order.  The
case of order zero is this line-sum statement.  For the induction step, the
global derivative identity just proved expresses every coordinate
directional derivative as a row with one additional raised slot; evaluate
that raised row at the line and use the induction hypothesis.  Since
coordinate directions span the configuration-space tangent space, all
multilinear positional jets vanish there.  This is the transverse-jet
condition used by the finite-scale argument.

First, finite test rows determine the free derivative data.  At each grade,
faithfulness of the free table says the family of all evaluation rows has
zero common kernel.  Since the field space at fixed grade is finite
dimensional, finitely many such rows already have zero common kernel: add
rows while the intersection of kernels drops in dimension; this can happen
only finitely often.  Under a scaling sequence, the prescribed local jet
convergence makes each of these finitely many actual rows converge on every
fixed free vector.  The limiting row map is injective.  Compose the actual
row maps with a left inverse of the limit; in finite dimension the resulting
operators converge to the identity and are invertible eventually.  If the
row maps are evaluated on the moving scaled derivative, the given fitted
row convergence then forces that derivative to converge to the free
derivative.  This is the gradient component of the first scaling bridge.

Second, the same finite-row method is applied to a finite family of
observers at adjacent grades.  The stress meter and the shell/rotation
flux identities identify the limiting rows with the free stress gradient
and free spin action.  Fixed-window jet convergence supplies convergence
of the actual rows.  The observer set is united before one scale threshold
is chosen, so these limits use one selected scale sequence.  The precise
rank-one argument identifying the positional scalar is given after the
finite collinear rows have been transported.

Third, transport the finite collinear tests through the same scale.  For
each selected scene $s$ and substitution $q$, the exact scaling identity
is
\[
 \mathrm{test}_k\!\left(s,\mathrm{aisle}_k(q)\right)
   =r_k^{m_s}\,\mathrm{lens}_{\mathrm{scaled}(k,s)}(q),
\]
where $m_s$ is the number of positional derivatives in that scene and
$r_k>0$ is the scale factor.  The local jet estimates make these row maps
converge on every fixed free substitution, and a finite-dimensional basis
then gives operator convergence for the finite family.  The derivative and
the local response are both transverse: each transported test row
annihilates each of them exactly at every scale.  Thus the same selected
family has the required common zero-row property; there is no injectivity
claim for the free collinear tests.  This is an exact-kernel argument, not
an eventual-injectivity argument.

Theorems~\ref{thm:free-quiet-classification} and
\ref{thm:free-collinear-scenes} give one more rank statement needed to
identify the normalization.  Work at each fixed grade $d$ above the common detector
threshold, where the free derivative $\partial$ is nonzero.  Let $V_d$ be
the finite-dimensional real space of admissible spin-sector response vectors
and let $R:V_d\to W_d$ be the free finite collinear row map.  Its kernel is
exactly the span of $\partial$.  Choose a real linear functional $\ell$ with
$\ell(\partial)\ne0$; then the augmented map $(R,\ell)$ is injective.
The selected row maps converge on a fixed basis, so adjoining this same
$\ell$ leaves them injective for all sufficiently large scales.  Choose one
such scale $\nu$.  Write $p_\nu$ and $\partial_\nu$ for the transported
response vectors, and let $A_\nu$ be the common aisle map from the original
finite-dimensional field space into the selected coordinates.  The source
row identity is $R_\nu A_\nu p=R_\nu A_\nu\partial=0$.  Since
$A_\nu$ is injective and $\partial\ne0$, $\partial_\nu\ne0$.  If
$\ell(\partial_\nu)=0$, then augmented injectivity of
$(R_\nu,\ell)$ would force $\partial_\nu=0$; thus this denominator is
nonzero.  Set $c_{d,\nu}=\ell(p_\nu)/\ell(\partial_\nu)$.  The
vector $p_\nu-c_{d,\nu}\partial_\nu$ is killed by both $R_\nu$ and
$\ell$, hence is zero.  Cancel the injective common aisle map to obtain
$p=c_{d,\nu}\partial$ in the original field space.  Since $p\ne0$,
$c_{d,\nu}\ne0$.  Grade inclusion intertwines $p$ and $\partial$.
Comparing adjacent grades above the detector threshold shows that this
scalar is constant there; comparing any lower grade with the fixed threshold
grade and using injectivity of inclusion gives the same identity below the
threshold.  For any finite collection of caps and grades, take the maximum
of the finitely many row-convergence thresholds.  Thus the same scalar $c$
works for every fixed finite comparison in the selected family.  The independent
stress-normalization result, together with the A05 observer limits,
identifies the spin response with the corresponding same-scale multiple;
define $\widetilde Y=c^{-1}Y$.  This ratio is obtained before the downstream
rotation-action argument.

Finally, erasure transfers from the free table to the selected table by
the same argument.  At fixed grade, regard the internal Ward maps as a
finite-dimensional real vector space.  Free erasure says the totality of
all separated row sums separates this space, so finitely many rows suffice.
Their actual counterparts converge on every fixed map.  Eventual
injectivity follows by applying a fixed left inverse of the free row map.
For a selected-table Ward map, pull it back by the aligned grade
isomorphism.  It lies in the free covariance subspace; if all selected row
sums vanish, the finite actual rows annihilate the pullback, hence it is
zero.  Conjugating back proves erasure at all sufficiently large grades.
The finite-dimensional injectivity argument uses pointwise convergence on
a fixed basis and does not assume convergence on a moving map.

\begin{lemma}[Differentiated Ward identity and brackets]
\label{lem:ward-brackets}
Suppose Still holds at every separated configuration; suppose also the
spin rows satisfy Euclidean covariance and reality, preserve the unit,
and commute with grade inclusion.  If the selected table has the erasure
property just described, then for every grade $d$ and $u,v\in\mathfrak{so}(4)$,
\[
 [\mathcal Y_d(u),\mathcal Y_d(v)] = \mathcal Y_d([u,v]).
\]
\end{lemma}

\begin{proof}
For a fixed word $q$ define the simultaneous orbital vector field
\[
 X_u(x)_i=u x_i.
\]
The Ward identity and the derivative rule give
\[
 X_u S(q;x)=-\sum_i S(q[i\mapsto\mathcal Y(u)q_i];x).
\tag{1}\label{eq:ward-orbital}
\]
The minus sign is fixed by the convention in the angular expression.
Because Still holds throughout the open configuration space, (1) can be
differentiated there and used again on each substituted word.  The vector
field bracket is $[X_u,X_v]=-X_{[u,v]}$; this follows directly by
differentiating the linear maps $x\mapsto ux$ and $x\mapsto vx$.
Consequently the commutator of the two orbital derivatives on a moment
is $-X_{[u,v]}S(q;x)$.

On the other hand, apply (1) twice.  Terms where the two spin substitutions
hit distinct slots cancel in pairs: the term indexed by $(i,j)$ in one
order is the same moment as the term indexed by $(j,i)$ in the reverse
order.  The terms with both substitutions at slot $i$ leave the positive
commutator
\[
 \sum_i S\bigl(q[i\mapsto
  ([\mathcal Y_d(u),\mathcal Y_d(v)]q_i)];x\bigr).
\]
The orbital vector-field bracket is $-X_{[u,v]}$, and by (1) its action on
the moment is $+\sum_i S(q[i\mapsto\mathcal Y_d([u,v])q_i];x)$.  Equating
the two computations gives zero row sum for the defect
\[
 D_d(u,v)=[\mathcal Y_d(u),\mathcal Y_d(v)]-\mathcal Y_d([u,v])
\]
has zero row sum on every separated word.

The covariance assumptions imply that $D_d$ is an internal Ward map:
it transforms under signed rotations, respects conjugation, annihilates
the unit, is alternating, and is compatible with grade inclusion.  Thus
it belongs to the finite-level internal kernel tested by erasure.  Erasure
annihilates it for every grade above a fixed threshold.  For any smaller
grade, raise to a grade above that threshold; inclusion compatibility
identifies the raised defect with the image of the original one.  The
grade inclusion is injective, so the original defect was zero as well.
This proves the bracket identity at every grade.
\end{proof}

\begin{corollary}[Selected finite-level Ward package]
\label{cor:selected-ward-package}
Assume the selected finite-level data supplied by
Section~\ref{sec:actual-stress-normalization} and the finite-row
comparisons above.  In particular, for one $c\ne0$ the
raw responses obey $p=cV$ and $Y=c\widetilde Y$, where $V$ is the genuine
positional derivative and $\widetilde Y$ is the normalized spin action.
Assume also the two free inputs: a finite collinear classifier for the
quiet derivative rows and erasure of any alternating two-generator internal
map whose separated Ward rows vanish, with the hypotheses used in the
finite-row comparisons above. Then the
actual collinear rotational identity and line identity hold, the Ward
identity holds at every separated configuration, and the spin rows satisfy
the $\mathfrak{so}(4)$ bracket relations.
\end{corollary}

\begin{proof}
The compact stress packet gives the raw collinear identity
$\mathcal A_K(p,Y)=0$.  By the common scalar normalization,
$\mathcal A_K(p,Y)=c\mathcal A_K(V,\widetilde Y)$, so the genuine derivative
and normalized spin action satisfy the collinear hypothesis of
Theorem~\ref{thm:ward-global-propagation}.  That theorem extends the
rotational identity to all of $\operatorname{Conf}_n$.  The finite-row bridges produce
the selected scale data and transfer erasure.  Apply
Lemma~\ref{lem:ward-brackets} to that same selected family.  Each use is
for a fixed grade/configuration; the finite maxima are taken before the
scale threshold, so the argument does not interchange a pointwise limit
with an uncountable family of tests.
\end{proof}

The cove, line, mass, point and exterior prices used by the stress packet
are the exact inputs of Section~\ref{sec:physical-ward-packet-prices}; their
exponent composition and the $-101+100=-1$ normalization are proved above.
The common selected curve, finite observer/jet convergence and scalar
normalization are supplied by Section~\ref{sec:actual-stress-normalization}.
For the Maxwell reference table, Theorems~\ref{thm:free-quiet-classification}
and~\ref{thm:free-collinear-scenes} supply the quiet-charge classifier and its
finite collinear tests at grades above their common threshold.
Theorem~\ref{thm:free-internal-erasure} supplies the internal-map erasure
input at the stated covariance, reality, unit and filtration hypotheses.
Each fixed grade and finite observer family can use its own comparison scale;
the scalar is common across grades by inclusion and nonvanishing, as proved
above. No scale threshold uniform over all grades is needed.
These free inputs are distinct from the packet-price bounds and from the
geometric line-to-configuration argument.
